% ==========================================================================
% GENERATED by scripts/split_paper.py. Do not edit this file.
%
% Source: imported_papers/sds_hounsfield_features/main.tex
% Part 4: Experimental Setup
%
% Edits belong upstream, in the published paper. To change what the
% chapter includes, change imports/hounsfield.toml and re-run
% the splitter.
% ==========================================================================
\section{Experimental Setup}

We evaluate our proposed approach using widely adopted deep learning architectures for medical image segmentation, training them with default hyperparameter settings. To ensure a rigorous assessment, we employ standard validation procedures in accordance with best practices.

\subsection{Model}
For the segmentation task, we build upon a 2D U-Net model~\cite{navab_u-net_2015}, a widely recognized state-of-the-art architecture in this domain, known for its effectiveness even in data-limited scenarios. The U-Net model follows an encoder-decoder structure with skip connections to preserve spatial information. To test our approach on different model sizes, we train three model variants with different initial feature map sizes ($16$, $32$, or $64$ channels). All configurations are based on a standard U-Net architecture and maintain a consistent depth of five encoding and decoding stages. They follow the standard practice of doubling the number of feature maps at each downsampling stage and ending with a $1\times1$ kernel for pixel-wise prediction \cite{navab_u-net_2015}. The models contain around 2M (\emph{U-Net-small}), 8M (\emph{U-Net-base}), and 31M (\emph{U-Net-large}) parameters, respectively.

Unlike the original U-Net architecture, we employ padded convolutions and add batch normalization after each convolutional layer in the encoder and decoder.
% No dropout or other regularization was employed.

\subsection{Data}
We conduct our experiments using CT data obtained from the Cantonal Hospital Aarau, Switzerland. 
Our dataset contains $176$ CT slices in $512 \times 512$ resolution from $44$ patients with malnutritional conditions.
We segment four tissue types: The skeletal muscle (SM), intramuscular adipose tissue (IMAT), visceral adipose tissue (VAT), and subcutaneous adipose tissue (SAT).
For each patient, we label four CT slices acquired at the vertebral levels: L1, L3, T10, and T12.
All patients gave general consent for research purposes, although due to data privacy regulations and national legislation (HRA), we are unable to publicly share imaging data, associated annotations, or model predictions.

To demonstrate the effectiveness of our approach on different dataset sizes, we sample subsets only containing fractions (\SI{10}{\percent}, \SI{20}{\percent}, $\dots$, \SI{100}{\percent}) of the available data.
The subsets are \emph{nested}, meaning a bigger subset includes smaller subsets, simulating the growth of an increasingly large dataset.

\subsection{Annotation process}
Following a thorough review, one case was removed due to unsuitable quality for manual analysis, as skeletal muscle tissue extended beyond the field of view. No further exclusion was needed due to, e.g., severe anatomical or technical abnormalities, such as body habitus variations, metal artifacts, or massive fluid accumulation. The remaining slices were classified based on characteristics commonly encountered in routine CT scans, including subcutaneous adipose tissue extending beyond the image field and mild image quality impairments (e.g., streaking or graininess). After dataset review, all eligible slices underwent semi-assisted manual segmentation using Slicomatic™ software from TomoVision at the specified vertebral levels \cite{baumgartner2023association}.

\subsection{Training parameters}
% No dedicated hyperparameter tuning was performed to isolate the impact of our feature engineering.
Each model is trained using the standard Adam optimizer~\cite{kingma2014adam} with its default settings: learning rate $\alpha=0.001$, $\beta_1=0.9$, $\beta_2=0.999$, and $\epsilon=10^{-8}$. The optimization process minimizes the cross-entropy loss function using a batch size of $4$.
Regardless of the dataset fraction used, we maintain a fixed number of \num{1800} optimization steps across all settings to ensure convergence and fair comparison of model performance.

\subsection{Validation}
All experiments are conducted using a 3-fold cross-validation across patients.
Folds are generated using the same random seed, enabling consistent comparisons across all experiments.

Performance metrics are reported as the mean value computed across the three folds. For each fold, all evaluation metrics are computed based on the best model checkpoint, as determined by the lowest validation loss.
We report the Dice score, which quantifies the overlap between the predicted segmentation ($P$) and the ground truth segmentation ($G$) and is defined as:
%
\begin{equation*}
\frac{2 \cdot |P \cap G|}{|P| + |G|}
\end{equation*}
%
Additionally, we report the precision, the recall, and the F$_1$ score.
All metrics are macro-averaged (calculated per tissue class and then averaged across classes), which is common practice to account for class imbalance.
To quantify the uncertainty in the estimated means, we report the \SI{95}{\percent} confidence interval, estimated using the standard deviation across folds. % Additionally, for each fold, all evaluation metrics are computed based on the best model checkpoint, as determined by the lowest validation loss.
